\section{Physical equation, provenance, and the pair balance}
\label{sec:context}

The parent manuscript \cite{BakerExterior2026} develops a complete
exterior-moment hierarchy and counted-block transport for the Newman
field. The base manuscript \cite{BakerVolterra2026} supplies its canonical
source and Jacobi--Volterra realization. We use their normalization
throughout:
\begin{align}
 r_k(u)&=\pi k^2e^{4u},&
 \Phi(u)&=\sum_{k\ge1} e^u\bigl(2r_k(u)^2-3r_k(u)\bigr)e^{-r_k(u)},
 \label{eq:source}\\
 H(t,z)&=\int_0^\infty e^{tu^2}\Phi(u)\cos(zu)\dd u,&
 H_t&=-H_{zz},\qquad H(0,z)=\tfrac18\xi((1+\ii z)/2).
 \label{eq:heat}
\end{align}
For $u\ge0$, the source is positive and decreases faster than
$e^{-Cu^2-Cu}$ for every fixed $C>0$. Compact differentiation of
\eqref{eq:heat} is therefore legitimate to every fixed order. For global
divisor statements we use the parent's established even order-one
canonical product, including its count, cutoff and complete
regularization hypotheses. The local block results below need only a
holomorphic heat solution and a fixed zero-free enclosing boundary.
They do not require an all-real exterior.

The founding exploration \cite{BakerConservationBrainstorm2026} contains
the finite identities and fourth-order block balance developed here.
The bounded extension in Sections~\ref{sec:block}--\ref{sec:errors}
selects a general homogeneous class and compares explicit error
allowances. The result is an identity and intermediate-estimate
investigation in the terminology of the project's result and avoidance
indices \cite{ProjectResults2026,ProjectAvoid2026}. A representation of
initial data does not by itself supply a bound on that data. A negative
result is always restricted to its displayed class. We make no priority
claim for finite heat-pole integrability: the abstract of
\cite{DeconinckKimuraSegur2007} reports conserved quantities for rational
Burgers poles, and the checked passages of \cite{Sokal2023HeatZeros}
discuss finite heat-zero dynamics and the specialized matrix in
Section~\ref{sec:calogero}. All formulas used here are derived below;
no uninspected external theorem or conjecture is imported.

\subsection{Count and pair data}
Put $B=H_z/H$. Off the divisor,
\begin{equation}
 B_t=-\partial_z(B_z+B^2).
 \label{eq:burgers}
\end{equation}
If one stationary positively oriented contour $\Gamma$ remains zero-free,
its integer count
\begin{equation}
 N=\av{B(t,z)}
 \label{eq:count}
\end{equation}
is constant. This follows either from continuity and the argument
principle, or by integrating the total derivative in
\eqref{eq:burgers}. Two simple zeros becoming one double zero do not
change the multiplicity count.

For a counted pair write its roots as $c\pm s$, set $D=s^2$, and let
$w_j=\rho_j-c$ for each exterior occurrence. The parent's complete
collar assumptions make the following sums absolutely convergent:
\begin{equation}
 S_n=\sum_j(w_j^2-D)^{-n}\quad(n\ge1),\qquad
 T_n=\sum_jw_j(w_j^2-D)^{-n}\quad(n\ge2).
 \label{eq:moments}
\end{equation}
In particular, $T_1$ is not part of this convergent coordinate family.
Writing $M=S_1$ and $\F=12S_2+16DS_3$, the regular pair equations are
\begin{equation}
 D'=2-4DM,\qquad M'=2M^2-\F.
 \label{eq:pair}
\end{equation}
The first law follows by subtracting the two simple-root velocities;
the mutual force contributes $2$ to $D'$, and each exterior root
contributes $-4D/(w_j^2-D)$. The second is the $n=1$ row of the
integer hierarchy displayed in Section~\ref{sec:obstruction}. The
parent's counted-contour passage extends both laws through mergers.

\begin{proposition}[An energy balance, not a general first integral]
Define
\begin{equation}
 E=2M-2DM^2+D\F.
 \label{eq:pair-energy}
\end{equation}
Then $E'=D\F'$. Also,
\begin{equation}
 (D')^2-2DD''+8\F D^2=4.
 \label{eq:pair-identity}
\end{equation}
For a prescribed forcing $\F(t)$, \eqref{eq:pair} is the
time-dependent Hamiltonian system $D'=E_M$, $M'=-E_D$, where partial
derivatives hold $\F(t)$ fixed.
\end{proposition}
\begin{proof}
Substitute \eqref{eq:pair} into
\[
 E'=2M'-2D'M^2-4DMM'+D'\F+D\F'.
\]
All terms except $D\F'$ cancel. Differentiating the first equation of
\eqref{eq:pair} gives $D''=-8M+8DM^2+4D\F$, whose substitution
proves \eqref{eq:pair-identity}. The two partial derivatives of
\eqref{eq:pair-energy} give the asserted Hamilton equations.
\end{proof}

The physical hierarchy determines $\F'$ through higher moments, so
constant forcing cannot be assumed. Subtracting $\int D\F'\dd t$
formally produces a constant but still requires that entire flux.
At a pair merger $D=0$, \eqref{eq:pair} gives $D'=2$, consistently
with \eqref{eq:pair-identity}. On a chart where $D>0$, putting
$r=\sqrt D$ gives
\begin{equation}
 r''=2\F r-r^{-3}.
\end{equation}
The inverse-cubic term is attractive. These formulas do not produce
a positive barrier against a merger; the $r$ chart itself ends there.
